\documentclass[11pt]{article}
\usepackage[T1]{fontenc}
\usepackage{lmodern,microtype,amsmath,amssymb,amsthm,mathtools,mathrsfs,tikz-cd}
\usepackage[margin=1in]{geometry}
\usepackage{xurl,hyperref}
\hypersetup{colorlinks=true,linkcolor=blue,citecolor=blue,urlcolor=blue,
pdftitle={Model interpretation of syntactic recursion categories},pdfauthor={Florian Lengyel}}
\newtheorem{theorem}{Theorem}[section]
\newtheorem{proposition}[theorem]{Proposition}
\newtheorem{lemma}[theorem]{Lemma}
\newtheorem{corollary}[theorem]{Corollary}
\newcommand{\N}{\mathbb N}
\newcommand{\PA}{\mathrm{PA}}
\newcommand{\Pfn}{\mathbf{Pfn}}
\newcommand{\ran}{\operatorname{ran}}
\newcommand{\Dom}{\operatorname{Dom}}
\title{Model interpretation of syntactic recursion categories}
\author{Florian Lengyel}
\date{30 October 2026}
\begin{document}
\maketitle

\begin{abstract}
For a consistent extension $T$ of Peano arithmetic, interpretation in
countable models gives a faithful range functor from the syntactic
recursion category $S_T$ to a locally connected split restriction category.
For computably enumerable $T$, Di Paola--Montagna's theorem says that
$S_T$ itself is not locally connected.
Theory extension is compatible with restriction of the family of models.
Every target switch associated with a complemented source domain is the
image of a single $T$-provably total program. The local-connectedness
inseparability theorem therefore applies to $S_T$ without a soundness
assumption. For computably enumerable $T$, numeral instances of domain
predicates give the separate effective comparison with classical recursive
inseparability and essential incompleteness.
\end{abstract}

\section{The syntactic recursion category}

Composition is written $gf=g\circ f$. In a restriction category, the
domain of $f:X\rightharpoonup Y$ is a restriction idempotent
$\bar f:X\rightharpoonup X$. The arrow is total when $\bar f=1_X$.
Domains are ordered by $a\le b$ when $ba=a$; their meet is composition.
For a domain $a$ on $Y$, write $f^*a=\overline{af}$. A domain $r$ on
$Y$ \emph{receives} $f$ when $rf=f$. A range of $f$
is the least domain receiving $f$, with the additional agreement property
\[
 gf=hf\quad\Longrightarrow\quad
 g\ran(f)=h\ran(f).
\]
We use restriction zeros and binary restriction coproducts, whose injections
are total. For $X+X$ these injections are denoted by $i_0,i_1$, and
$\nabla_X=[1_X,1_X]:X+X\to X$ denotes the codiagonal.
Complementary domains $c,c'$ satisfy $cc'=0$ and $c\vee c'=1$.
A disjoint pair $a_0,a_1$ is \emph{inseparable} if no complementary pair
satisfies $a_0\le c$ and $a_1\le c'$.

Let $T$ be a consistent classical theory in the language of arithmetic
containing $\PA$. Formal numerical types are generated from $\omega$ by
binary products and tagged coproducts. Fix recursive numerical codings
whose bijection and structural equations are provable in $\PA$.
For a typed program index $e:X\rightharpoonup Y$, let $G_e(x,y)$ be its
$\Sigma_1$ output formula. It existentially quantifies a computation trace
over a primitive-recursive trace test. Fix a deterministic programming
system whose graph identities, including functionality, composition,
and parameterization, are provable in $\PA$.

The category $S_T$ has these formal types as objects and typed indices
modulo the equivalence relation
\begin{equation}\label{eq:uniform}
 e\equiv_T f
 \quad\Longleftrightarrow\quad
 T\vdash \forall x\forall y\,
       \bigl(G_e(x,y)\leftrightarrow G_f(x,y)\bigr).
\end{equation}
The quantifiers belong to the single sentence proved in $T$.
This is the unprimed syntactic construction of Montagna and
Di Paola--Montagna \cite{Montagna,DPM}.

For a program operation $\operatorname{comp}$ implementing composition,
$\PA$ proves
\begin{equation}\label{eq:composition}
 G_{\operatorname{comp}(f,e)}(x,z)
 \leftrightarrow
 \exists y\,\bigl(G_e(x,y)\land G_f(y,z)\bigr).
\end{equation}
Substitution of uniformly equivalent graphs under this existential proves
that composition descends to \eqref{eq:uniform}. The identity and
associativity equations follow from the same graph identities.

Every $\Sigma_1$ formula $C(x)$ represents a partial identity: search for
a witness to $C(x)$ and then return $x$. Write $[C]_T$ for this domain.
Then
\begin{equation}\label{eq:domain-order}
 [C]_T\le[D]_T
 \quad\Longleftrightarrow\quad
 T\vdash\forall x\,(C(x)\longrightarrow D(x)).
\end{equation}
Conjunction and disjunction give domain meets and joins. The domain and
range of $[e]_T$ are represented respectively by
\[
 \exists y\,G_e(x,y),\qquad
 R_e(y):=\exists x\,G_e(x,y).
\]
For the range property, a domain $[C]_T$ receives $[e]_T$ precisely when
$T$ proves $G_e(x,y)\to C(y)$ universally, equivalently when
$T$ proves $R_e(y)\to C(y)$ universally. For agreement, a witness to
$R_e(y)$ together with equality of two composite graphs gives equality
of the two output graphs at $y$. Existential elimination in $T$ proves
the required agreement equation. These arguments also show that the
chosen ranges are unchanged under theory extension.

The ordinary product and tagged-coproduct programs supply the restriction
products and coproducts. Write $\pi_1,\pi_2$ for the projections and
$\Delta_X=\langle1_X,1_X\rangle:X\to X\times X$ for the total categorical
diagonal. The universal program supplies a Turing morphism
\[
 \tau:\omega\times\omega\rightharpoonup\omega:
 \quad f:X\times\omega\rightharpoonup\omega
 \ \Longrightarrow\ 
 f=\tau(r\times1_\omega)
\]
for some total parameter arrow $r:X\to\omega$. Formalized
parameterization makes this arrow $\PA$-provably total. The chosen coding
isomorphisms make $S_T$ isotypical, and product with a zero arrow is zero.
Thus $S_T$ is a prodominical recursion category. The preceding
constructions do not require that the axioms of $T$ be computably enumerable.

\section{Interpretation in models}

The category $\Pfn$ has sets as objects and partial functions as arrows,
with ordinary composition. Restrictions and ranges are partial identities
on domains of definition and images. Restriction products are Cartesian products;
coproducts are tagged disjoint unions.

Fix a model $M\models T$. Interpret numerical types recursively by
\[
 X_M=|M|\quad(X=\omega),\qquad
 (X\times Y)_M=X_M\times Y_M,\qquad
 (X+Y)_M=X_M\amalg Y_M.
\]
The fixed arithmetic codings identify these sets with the corresponding
codes in $M$. A standard program index in a graph formula is interpreted
as its numeral in $M$. Define
\begin{equation}\label{eq:model-functor}
 F_M(X)=X_M,\qquad
 F_M([e]_T)(x)=y
 \quad\Longleftrightarrow\quad M\models G_e(x,y).
\end{equation}
All quantified traces and intermediate values in this formula range over
the whole model, including its nonstandard elements.

\begin{proposition}\label{prop:model}
For every $M\models T$, \eqref{eq:model-functor} defines a covariant
functor $F_M:S_T\to\Pfn$ preserving restrictions, zeros, binary
restriction products and coproducts, ranges, and finite joins of domains.
\end{proposition}
\begin{proof}
Functionality is provable in $\PA$, so the interpreted graph is a
partial function. Equation \eqref{eq:uniform} makes the assignment
independent of the representative. Equation \eqref{eq:composition},
interpreted in $M$, is exactly relational composition of the two partial
functions; their intermediate values range over $Y_M$. The identity
graph gives the identity function.

The domain predicate $\exists y\,G_e(x,y)$ gives the interpreted
function's domain. Similarly,
\[
 \{y\in Y_M:M\models R_e(y)\}
 =\operatorname{im}(F_M[e]_T),
\]
so ranges are preserved. False predicates, conjunction, and disjunction
give zeros, intersections, and unions. The PA-provable product and
tagged-coproduct graph equations give the corresponding operations on
partial functions.
\end{proof}

\section{A faithful family and theory extension}

Let $\mathcal M_T$ be the set of all structures for the language of
arithmetic with underlying set $\N$ that satisfy $T$. Their arithmetic
operations need not be the standard ones, and the interpretation of a
numeral $\bar n$ need not be the carrier element $n$. Completeness and
the downward L\"owenheim--Skolem theorem imply
$\mathcal M_T\ne\varnothing$: every model of $T$ is infinite, and a
countably infinite model can be transported to the carrier $\N$.

Define the product category
\[
 \mathcal D_T=\prod_{M\in\mathcal M_T}\Pfn.
\]
An object is a family of sets $(A_M)_M$; an arrow is a family of partial
functions with the indicated sources and targets. Composition and the
restriction, range, restriction-product, and coproduct operations are coordinatewise.
The functors of Proposition~\ref{prop:model} combine to give
\[
 F_T:S_T\longrightarrow\mathcal D_T,\qquad
 F_T(X)=(X_M)_M,\quad F_T([e]_T)=(F_M[e]_T)_M.
\]

A domain $a$ is \emph{connected} if $a=b\vee c$ and $bc=0$ imply
$b=0$ or $c=0$. A category is \emph{locally connected} when its connected
domains are jointly epimorphic: agreement of parallel arrows after every
connected restriction implies equality. A category is \emph{split} when
every restriction idempotent splits. A partial monomorphism
$p:X\rightharpoonup Y$ satisfies
$pu=pv\Rightarrow\bar p u=\bar p v$ for all $u,v:W\rightharpoonup X$.
\emph{Weak choice} means that
every partial monomorphism has a section $s$ with
$ps=\bar s$ and $psp=p$.

\begin{theorem}\label{thm:faithful}
The functor $F_T$ is faithful and preserves restrictions, zeros, binary
restriction products and coproducts, ranges, and finite domain joins. The category
$\mathcal D_T$ is split and locally connected, satisfies weak choice,
and has inverse images preserving binary domain joins. Moreover,
$F_T(1_\omega)\ne0$.
\end{theorem}
\begin{proof}
Preservation follows coordinatewise from Proposition~\ref{prop:model}.
If $[e]_T\ne[f]_T$, let $E$ denote the uniform graph-equality sentence
in \eqref{eq:uniform}. The theory $T\cup\{\neg E\}$ is consistent.
Completeness and downward L\"owenheim--Skolem give a model on carrier
$\N$ in which the two graphs differ. This is a coordinate of $F_T$,
which therefore distinguishes the arrows.

A domain in $\mathcal D_T$ is a family of subsets. Its splitting uses
those subsets as the coordinate sets. For an object $A=(A_M)_M$, a
domain supported at one coordinate $M$, with a singleton subset
$\{x\}\subseteq A_M$ there, is connected. These domains detect the
definedness and value of every coordinate of a partial arrow. Hence
$\mathcal D_T$ is locally connected.

A partial monomorphism has injective coordinate functions: test its
defining property on singleton objects supported at one coordinate.
The partial inverses of these injections form a section, proving weak
choice. Inverse images commute with unions coordinatewise. Finally,
$\mathcal M_T$ is nonempty and every interpreted numerical object is
nonempty, so its identity differs from zero.
\end{proof}

For consistent computably enumerable $T\supseteq\PA$, the source $S_T$
is not locally connected \cite[Corollary 4.4, pp.~657--658]{DPM}.
The connected domains used in the preceding proof belong to
$\mathcal D_T$ and need not be images of source domains. Faithfulness
reflects equality of arrows; it does not require these additional target
domains to have source preimages.

For an inclusion $T\subseteq T'$ of consistent theories, define
\[
 Q_{T,T'}:S_T\to S_{T'},\qquad
 Q_{T,T'}X=X,\quad Q_{T,T'}[e]_T=[e]_{T'}.
\]
Provability is monotone, so this is a well-defined covariant functor.
It preserves the chosen computational operations and Turing morphism.
Since $\mathcal M_{T'}\subseteq\mathcal M_T$, restriction of families
defines a covariant projection
$\pi_{T,T'}:\mathcal D_T\to\mathcal D_{T'}$. Fixing type interpretations
once for each arithmetic structure gives a strictly commuting square:
\[
\begin{tikzcd}[column sep=large]
 S_T \arrow[r,"Q_{T,T'}"] \arrow[d,"F_T"'] &
 S_{T'} \arrow[d,"F_{T'}"]\\
 \mathcal D_T \arrow[r,"\pi_{T,T'}"'] & \mathcal D_{T'}.
\end{tikzcd}
\]
Both paths send a program to the family of its interpretations in the
models of $T'$. The identity and composition laws for the projections
and quotient functors follow from restriction of indices and the
unchanged program representatives.

\section{Switching and the local-connectedness theorem}

A \emph{range functor} preserves restrictions, ranges, restriction zeros,
and binary restriction coproducts. Since $c\vee c'=\ran[c,c']$,
such a functor preserves binary domain joins and complementary pairs.
Theorem~\ref{thm:faithful} supplies
such a functor into a split locally connected category satisfying weak
choice and inverse-image preservation of domain joins. These are the
target hypotheses of the local-connectedness inseparability theorem
\cite[Theorem 4.7]{LOC}. Its remaining requirement concerns lifting
specific sections associated with complemented source domains.

\subsection{The local-connectedness construction}

The following argument supplies the needed sections. It is the
construction of \cite[Section 4.3]{LOC}, with its restriction-category
hypotheses stated explicitly. For a coproduct $Y+Z$, let
$\varepsilon_0=1_Y+0_Z$ and $\varepsilon_1=0_Y+1_Z$ be its summand
domains. Their meet is zero and their join is the identity. These facts
follow by precomposition with the two injections and the agreement
property of ranges.

\begin{lemma}\label{lem:connected-alternative}
Let $\mathcal A$ be a restriction category with ranges, restriction zeros,
and binary restriction coproducts. Suppose inverse images preserve binary
domain joins. If $a$ is a connected domain on $X$ and
$f:X\to Y+Z$ is total, then
$\varepsilon_0fa=fa$ or $\varepsilon_1fa=fa$.
\end{lemma}
\begin{proof}
Totality and the inverse-image laws give the disjoint decomposition
\[
 a=af^*\varepsilon_0\vee af^*\varepsilon_1.
\]
Connectedness makes one term zero and the other equal to $a$. In the
latter case $a\le f^*\varepsilon_i$, and the restriction identity
$\varepsilon_i f=f\overline{\varepsilon_i f}$ gives
$\varepsilon_i fa=fa$.
\end{proof}

\begin{proposition}\label{prop:local-switch}
Let $\mathcal A$ be a locally connected split restriction category with
ranges, restriction zeros, and binary restriction coproducts. Suppose
weak choice holds and inverse images preserve binary domain joins.
For disjoint domains $a,b$ on $X$, the arrow
$p=\nabla_X(a+b):X+X\rightharpoonup X$ is a partial monomorphism
with a unique section $s:X\rightharpoonup X+X$, and
\[
 ps=\bar s=a\vee b.
\]
In particular, for complementary $c,c'$ on $X$, the arrow
$\nabla_X(c'+c)$ has a total section $s:X\to X+X$ satisfying
$sc=i_1c$ and $sc'=i_0c'$.
\end{proposition}
\begin{proof}
First let $\theta,\theta':Y\to X+X$ be total and suppose
$p\theta=p\theta'$. The identities
$ap=\nabla_X(a+0)$ and $bp=\nabla_X(0+b)$ give
\[
 \nabla_X(a+0)\theta=\nabla_X(a+0)\theta',
 \qquad
 \nabla_X(0+b)\theta=\nabla_X(0+b)\theta'.
\]
For a connected domain $\gamma$ on $Y$,
Lemma~\ref{lem:connected-alternative} places each of
$\theta\gamma,\theta'\gamma$ in one summand.
If both lie in the first summand, the first equality, followed by $i_0$,
gives $(a+b)\theta\gamma=(a+b)\theta'\gamma$; the second-summand
case uses the second equality and $i_1$.
If they lie in opposite summands, the two equalities force both these
restricted arrows to be zero. Local connectedness therefore gives
$(a+b)\theta=(a+b)\theta'$. Since $\bar p=a+b$, this proves the
partial-monomorphism condition for total test arrows.

For arbitrary $u,v:W\rightharpoonup X+X$ with $pu=pv$, split their
common composite domain as
$\overline{pu}=\overline{pv}=me$, where $m:Z\to W$ is total and
$em=1_Z$. The arrows $pum,pvm$, hence also $um,vm$, are total.
The preceding case gives $\bar p um=\bar p vm$. Composing with $e$
and using $\bar p u=\bar p u\,\overline{pu}$ and its analogue for $v$
gives $\bar p u=\bar p v$. Thus $p$ is a partial monomorphism.

Weak choice supplies a section $s$. The equations $ps=\bar s$ and
$psp=p$ imply $\bar s=\ran(p)$: this domain receives $p$, and every
domain receiving $p$ also receives $ps=\bar s$. Moreover
$\bar p s=s\overline{ps}=s$. Two sections have the same composite
with $p$, so partial monicity makes them equal.
The least-domain property gives $\ran(p)=a\vee b$.

For $a=c'$ and $b=c$, the section is total because $c'\vee c=1_X$.
The equality $p(sc)=c=p(i_1c)$ and partial monicity give $sc=i_1c$;
the equation on $c'$ is analogous.
\end{proof}

\begin{theorem}[Range-functor lifting]\label{thm:local-lifting}
Let $\mathscr R$ be a prodominical recursion category with ranges and
binary restriction coproducts. Let $\mathcal A$ be a locally connected
split restriction category with ranges, restriction zeros, and binary
restriction coproducts, satisfying weak choice, with inverse images
preserving binary domain joins. Let
$F:\mathscr R\to\mathcal A$ be a range functor.
Choose a Turing morphism $\tau:U\times U\rightharpoonup U$ and
an isomorphism $b:U\to U+U$, with $F(1_U)\ne0$.
For each complementary pair $c,c'$ on $U$, let $s_c$ be the section
of $\nabla_{FU}(Fc'+Fc)$. Suppose there is an arrow
$\ell_c:U\rightharpoonup U+U$ with $F\ell_c=s_c$, using the coproduct
comparison. Then
\[
 a_0=(b\tau\Delta_U)^*(1_U+0_U),\qquad
 a_1=(b\tau\Delta_U)^*(0_U+1_U)
\]
are inseparable in $\mathscr R$.
\end{theorem}
\begin{proof}
This is \cite[Theorem 4.7]{LOC}; the following proof uses the
inverse-image laws directly after constructing the section.
The domains $a_0,a_1$ are disjoint. Suppose $c,c'$ separate them.
Put $u=b^{-1}\ell_c$, and choose a total parameter $g:U\to U$ with
$\tau(g\times1_U)=u\pi_2$. Precomposition with $\langle1_U,g\rangle$
gives $\tau\Delta_Ug=ug$.

In the target write $E=Fg$, $H=F(b\tau\Delta_U)$, $\Gamma=Fc$,
and $\Gamma'=Fc'$. The preceding equation gives $HE=s_cE$.
Since $s_c\Gamma=i_1\Gamma$ and $s_c\Gamma'=i_0\Gamma'$,
its inverse images of the summand domains are $\Gamma'$ and $\Gamma$,
respectively. Preservation of restrictions and composition now gives
\[
 E^*F(a_0)=E^*\Gamma',\qquad E^*F(a_1)=E^*\Gamma.
\]
The separator inequalities imply that $E^*\Gamma$ and $E^*\Gamma'$
are equal. They are complementary because $E$ is total and target
inverse images preserve joins. Hence $1_{FU}=0$, contradicting the
hypothesis. The proof uses totality of $F\ell_c=s_c$, not totality of
its source lift $\ell_c$.
\end{proof}

\subsection{Lifting in the syntactic category}

\begin{lemma}\label{lem:switch}
Let $c,c'$ be complementary domains on $X$ in $S_T$. There is a total
arrow $s:X\to X+X$ with
\[
 sc=i_1c,\qquad sc'=i_0c'.
\]
Its image under $F_T$ is the unique section of
$\nabla_{F_TX}(F_Tc'+F_Tc)$, using the coproduct comparison to identify
$F_T(X+X)$ with $F_TX+F_TX$.
\end{lemma}
\begin{proof}
Choose $\Sigma_1$ representatives $C(x)=\exists u\,R(x,u)$ and
$C'(x)=\exists v\,R'(x,v)$ with primitive-recursive witness tests.
The theory $T$ proves universally their disjointness and coverage.
Search increasing finite bounds for a witness to either predicate,
returning the tagged input $i_1x$ on finding a $C$-witness and $i_0x$
on finding a $C'$-witness. Fix a priority for simultaneous discoveries.
The finite-search verification in $\PA$ proves that the program halts
exactly when $C(x)\lor C'(x)$. Under the disjointness assumption it also
proves the exact graph equation
\[
 G_s(x,z)\leftrightarrow
 (C(x)\land z=i_1x)\lor(C'(x)\land z=i_0x).
\]
Thus $T$ proves this equation and universal halting. The equations for
the restrictions follow, giving a total arrow of the uniform quotient.

In each $M\models T$, the interpretations of $C,C'$ partition $X_M$.
The interpreted program sends each input to its specified tagged copy.
It is exactly the inverse of the partial injection
$\nabla_{X_M}(c'_M+c_M)$. Its graph is therefore the unique required
section. The same program works at every coordinate, proving the lifting
statement for the entire family.
\end{proof}

\begin{corollary}\label{cor:pair}
For every consistent $T\supseteq\PA$, the functor $F_T$ satisfies the
lifting hypothesis of the local-connectedness inseparability theorem.
Fix $U=\omega$, a Turing morphism $\tau:U\times U\rightharpoonup U$,
and a PA-provable coding isomorphism $b:U\to U+U$. Put
\[
 d=\tau\Delta_U,\qquad
 a_0=(bd)^*(1_U+0_U),\qquad a_1=(bd)^*(0_U+1_U).
\]
Then $a_0,a_1$ are inseparable in $S_T$. Their images under
$Q_{T,T'}$ are the identically defined inseparable pair in $S_{T'}$.
\end{corollary}
\begin{proof}
Theorem~\ref{thm:faithful} and Lemma~\ref{lem:switch} verify all the
hypotheses of Theorem~\ref{thm:local-lifting}. The quotient functor keeps
every program in the definition fixed, and the same argument applies
to consistent $T'$.
\end{proof}

The local-connectedness theorem uses Turing universality in its source.
The target $\mathcal D_T$ has all families of partial functions as arrows;
$F_T(\tau)$ is not universal for these arrows: at each model coordinate,
countably many program parameters cannot enumerate all partial functions
on its countable numerical carrier. The target's domain lattices are
Boolean, so the image of every disjoint source pair has a separator
there. Lemma~\ref{lem:switch} lifts the sections associated with source
complemented domains. An arbitrary target separator need not come from
such a domain. This distinction explains how a faithful range functor
can have an inseparable source pair and a separable image pair.

\section{Effective observations and essential incompleteness}

Let $\mathrm{Rec}$ be the classical recursion category on the formal
numerical types, with partial recursive functions as arrows and extensional
equality. Its domains are computably enumerable subsets, and its
complemented domains are recursive subsets.
Assume now that $T$ has computably enumerable axioms. For a domain
$c=[C]_T$ on $\omega$, define the subset of the standard natural numbers
\begin{equation}\label{eq:observation}
 E_T(c)=\{n\in\N:T\vdash C(\bar n)\}.
\end{equation}
It is computably enumerable by proof enumeration and is independent of
the representative. The operation is monotone, preserves finite meets,
and preserves top and bottom. Let $\kappa_n:\omega\to\omega$ be the
total constant program with value $n$. Then
\begin{equation}\label{eq:faithful-observation}
 n\in E_T(c)
 \quad\Longleftrightarrow\quad c\kappa_n=\kappa_n
 \quad\Longleftrightarrow\quad
 F_T(c)F_T(\kappa_n)=F_T(\kappa_n).
\end{equation}
The first equivalence follows from the partial-identity graph equation;
the second is faithfulness. Thus the standard observations can also be
expressed by equations in the faithful representation.

Take $b(2m)=i_0m$ and $b(2m+1)=i_1m$ in
Corollary~\ref{cor:pair}, using an acceptable universal program.
Let $\alpha(x),\beta(x)$ be the resulting $\Sigma_1$ domain formulas.
Their standard truth sets
\[
 A=\{n:\N\models\alpha(\bar n)\},\qquad
 B=\{n:\N\models\beta(\bar n)\}
\]
are the classical even-output and odd-output sets for self-application.
To prove inseparability, suppose a recursive set $C$ separated them.
Choose a program with index $e$ that returns $1$ on inputs in $C$
and $0$ on inputs outside $C$. If $e\in C$, its output on $e$ puts
$e$ in $B$, contrary to separation. If $e\notin C$, its output puts
$e$ in $A$, again a contradiction.
Each finite computation gives a PA-proof of its numeral instance, and $\PA$ proves
$\forall x\,\neg(\alpha(x)\land\beta(x))$.

Define
\[
 P_T=\{n:T\vdash\alpha(\bar n)\},\qquad
 N_T=\{n:T\vdash\neg\alpha(\bar n)\}.
\]
These are computably enumerable disjoint sets with $A\subseteq P_T$
and $B\subseteq N_T$. They enlarge under theory extension.
Let $\operatorname{Pair}(\mathrm{Rec})$ have objects $(X,A,B)$ consisting
of a numerical type and disjoint computably enumerable domains.
An arrow $(X,A,B)\to(Y,P,N)$ is a total recursive function $f:X\to Y$
with $f(A)\subseteq P$ and $f(B)\subseteq N$. Identities and composition
are inherited from $\mathrm{Rec}$. The assignment
$T\mapsto(\omega,P_T,N_T)$, with the identity function for theory
inclusions, is a covariant functor from the poset of consistent
computably enumerable extensions of $\PA$ to
$\operatorname{Pair}(\mathrm{Rec})$.
The inclusions of $A,B$ give identity-arrow reductions to every such pair.
A recursive separator pulls back along a total recursive function;
consequently $P_T,N_T$ are recursively inseparable.

\begin{corollary}
Every consistent computably enumerable extension of $\PA$ is incomplete.
\end{corollary}
\begin{proof}
If $T$ were complete, parallel enumeration of proofs of
$\alpha(\bar n)$ and $\neg\alpha(\bar n)$ would halt on every standard
input. Consistency makes the answers exclusive. The resulting recursive
set $P_T$ would separate $A$ and $B$, a contradiction.
\end{proof}

The effective step uses standard numeral instances and the presentation
of $T$ by computably enumerable axioms. The semantic functor exists also
for $T=\operatorname{Th}(\N)$, when $T$ is complete and $S_T$ is the
classical category of partial recursive functions on the formal numerical
types. Thus internal inseparability alone does not imply logical
incompleteness. Bernardi's precomplete-equivalence method gives another
established route through PA-provable program equivalence
\cite[pp.~125--128]{Bernardi}.

\section{The three-point alternative}

The same local-connectedness theorem applies to the computational cosheaf
category on the three-point space of \cite{ThreePoint}. Its basic diagram
shape is $J=(1\leftarrow0\to2)$. Let $\mathcal C_{\mathrm{rec}}$ have
computably enumerable numerical carriers and partial recursive functions
total on those carriers. The split ambient category is
\[
 \mathcal A=\operatorname{Par}
   (\mathcal C_{\mathrm{rec}}^J,\operatorname{Mono}),
\]
whose arrows are partial maps with monic domain inclusions. Let $X$ be
the universal diagram constructed in \cite{ThreePoint}, and let
$\mathscr R_X$ be the full isotypical recursion category on the formal
objects generated by $X$.
Its full inclusion $\mathscr R_X\hookrightarrow\mathcal A$ preserves
the required ranges and coproducts.

The subdiagram generated by a central element and its two endpoint
images is finite and connected. Singleton subdiagrams at either endpoint
are connected as well. All are computational domains. They detect every
component domain and value of a partial arrow, so $\mathcal A$ is locally
connected. Its splitness, weak choice, ranges, coproducts, and
componentwise inverse-image laws are established in \cite{ThreePoint}.
The section associated with a complemented domain has source and target
in $\mathscr R_X$; fullness therefore supplies its lift. This gives the
same switching argument in a nonnumerical prodominical recursion category.

\begin{thebibliography}{9}
\bibitem{Bernardi}
C.~Bernardi, \emph{Precomplete equivalence relations in dominical
categories}, Proc. Amer. Math. Soc. \textbf{98} (1986), 124--134.
\bibitem{DPM}
R.~A. Di Paola and F.~Montagna, \emph{Some properties of the syntactic
$P$-recursion categories generated by consistent, recursively enumerable
extensions of Peano arithmetic}, J. Symbolic Logic \textbf{56} (1991),
643--660.
\bibitem{LOC}
F.~Lengyel, \emph{Locally connected recursion categories}, manuscript,
revised 10 October 2026, \texttt{LOCOct10f.tex}.
\bibitem{ThreePoint}
F.~Lengyel, \emph{Range completion and computational cosheaves on a
three-point space}, 29 October 2026,
\href{https://github.com/flengyel/CategoricalEssentialIncompleteness/blob/3be65406b0bb6d3109a41fd3f517ca62fcc5ed8f/manuscript/range_completion_three_point_cosheaves.tex}{source}.
\bibitem{Montagna}
F.~Montagna, \emph{``Pathologies'' in two syntactic categories of partial
maps}, Notre Dame J. Formal Logic \textbf{30} (1989), 105--116.
\end{thebibliography}
\end{document}
